%%%%%%%%%%  scan idx 294  |  folio 281  %%%%%%%%%%
\item l'%
objet final de \(X\) est quasi-compact;
\item \(X\) a assez de points, et, pour tout point \(x\) de \(X\), tout \(A_x\)-module à gauche est de tor-dimension finie.
\end{enumerate}
Alors on a, d'après (I 5.10), \(D({}_A X)_{\parf}=D^b({}_A X)_{\coh}\), et l'homomorphisme (2.5.1) est un isomorphisme.

Les hypothèses (i) et (ii) sont satisfaites par exemple dans le cas d'un espace compact annelé en anneaux locaux réguliers.

\textbf{2.6. Accouplements entre \(K^{\bullet}\) et \(K^{\bullet}\) (cas non commutatif).} Soit
\[
 f:(X',A')\longrightarrow (X,A)
\]
un morphisme de topos annelés. Notons \(\Mod({}_{A'} X'_A)\) la catégorie des bi-Modules, à gauche sur \(A'\), à droite sur \(f^{-1}(A)\), et \(\Loclib_{A'}({}_{A'} X'_A)\) la sous-catégorie pleine formée des objets qui appartiennent à \(\Loclib({}_{A'} X')\). Par analogie avec (2.2), nous poserons
\[
 K^{\bullet}_{\naif}({}_{A'} X'_A)_{A'}=k(\Loclib_{A'}({}_{A'} X'_A)).
\]
Le produit tensoriel
\[
 \otimes_A:\Mod({}_{A'} X'_A)\times\Mod({}_A X)\longrightarrow \Mod({}_{A'} X')
\]
\[
 (E',E)\longmapsto E'\otimes_{f^{-1}(A)}f^{-1}(E)
\]
induit un foncteur
\begin{equation}\tag{2.6.1}
 \Loclib_{A'}({}_{A'} X'_A)\times\Loclib({}_A X)\longrightarrow\Loclib({}_{A'} X').
\end{equation}
L'application \((E',E)\mapsto\cl(E'\otimes_A E)\) déduite de (2.6.1) est biadditive au sens de (1.5), et définit par suite un homomorphisme
\begin{equation}\tag{2.6.2}
 K^{\bullet}_{\naif}({}_{A'} X'_A)_{A'}\otimes K^{\bullet}_{\naif}({}_A X)\longrightarrow K^{\bullet}_{\naif}({}_{A'} X').
\end{equation}
Notons \(D({}_{A'} X'_A)\) la catégorie dérivée de \(\Mod({}_{A'} X'_A)\), \(D({}_{A'} X'_A)_{A'\text{-}\parf}\) la sous-catégorie pleine formée des complexes d'amplitude parfaite finie rel. à \(A'\) (i.e. en tant qu'objets de \(D({}_{A'} X')\)), et posons, comme en (2.2),
\[
 K^{\bullet}({}_{A'} X'_A)_{A'}=k(D({}_{A'} X'_A)_{A'\text{-}\parf}.
\]
